
\documentclass[11pt]{article}
\usepackage[margin=1in]{geometry}
\usepackage{amsmath,amssymb,booktabs,array,longtable}
\usepackage{graphicx}
\usepackage{hyperref}
\usepackage{enumitem}
\usepackage{titlesec}

\begin{document}

\begin{center}
\Large\textbf{Shiv Nadar University Chennai}\\[2mm]
\end{center}

\renewcommand{\arraystretch}{1.4}

\begin{tabular}{|p{3.5cm}|p{8cm}|p{2cm}|}
\hline
Degree \& Branch & B.Tech Artificial Intelligence \& Data Science & Semester V\\ \hline
Subject Code \& Name & CS3807 -- Deep Learning Laboratory & AY: 2026--27\\ \hline
\end{tabular}

\vspace{3mm}

\begin{center}
\Large\textbf{Experiment 1}\\
\large\textbf{Implementation of a Single Layer Perceptron for Binary Classification}
\end{center}

%----------------------------------------------------------

\section*{1. Objective}

The objective of this experiment is to understand the concept of an artificial neuron and implement a Single Layer Perceptron from scratch for binary classification. Students will learn the perceptron learning algorithm, understand the importance of activation functions, visualize the learning process, and evaluate the classifier using a real-world dataset.

%----------------------------------------------------------

\section*{2. Background Theory}

An Artificial Neuron is the fundamental building block of a neural network. It receives one or more input features, multiplies them by corresponding weights, adds a bias term, and applies an activation function to determine the output.

The Single Layer Perceptron is the earliest supervised learning algorithm capable of solving linearly separable binary classification problems.

\subsection*{Mathematical Formulation}

Weighted Sum

\[
z=\mathbf{w}^{T}\mathbf{x}+b
\]

where

\begin{itemize}
\item $\mathbf{x}$ : Input vector
\item $\mathbf{w}$ : Weight vector
\item $b$ : Bias
\item $z$ : Linear combination
\end{itemize}

Prediction

\[
\hat y=f(z)
\]

Weight Update Rule

\[
\mathbf{w}_{new}
=
\mathbf{w}_{old}
+
\eta
(y-\hat y)
\mathbf{x}
\]

Bias Update

\[
b_{new}
=
b_{old}
+
\eta
(y-\hat y)
\]

where $\eta$ denotes the learning rate.

%----------------------------------------------------------

\section*{3. Activation Functions}

An activation function determines the output of a neuron based on the weighted sum of its inputs. It introduces non-linearity, enabling neural networks to learn complex patterns. Although the perceptron uses only the Step Activation Function, modern deep learning models employ differentiable activation functions for efficient optimization through backpropagation.

\subsection*{Step Activation Function}

\[
f(z)=
\begin{cases}
1,&z\ge0\\
0,&z<0
\end{cases}
\]

The Step Activation Function produces a binary output and is therefore suitable only for linearly separable binary classification problems.

\begin{center}
\begin{tabular}{|p{2.5cm}|p{3cm}|p{6cm}|}
\hline
\textbf{Activation} & \textbf{Output Range} & \textbf{Typical Applications}\\
\hline
Step & $\{0,1\}$ & Classical Perceptron\\
\hline
Sigmoid & $(0,1)$ & Binary Classification Output Layer\\
\hline
Tanh & $(-1,1)$ & Hidden Layers (Traditional Networks)\\
\hline
ReLU & $[0,\infty)$ & Hidden Layers in CNNs and MLPs\\
\hline
Leaky ReLU & $(-\infty,\infty)$ & Avoiding Dying ReLU\\
\hline
Softmax & $(0,1)$ & Multi-class Classification Output Layer\\
\hline
\end{tabular}
\end{center}

\textbf{Note:} This experiment uses only the Step Activation Function. The remaining activation functions will be studied in subsequent experiments involving Multilayer Perceptrons and Deep Neural Networks.

%----------------------------------------------------------

\section*{4. Dataset Description}

\textbf{Dataset:} Banknote Authentication Dataset

\textbf{Source:} UCI Machine Learning Repository

\textbf{Download Link:}

\url{https://archive.ics.uci.edu/dataset/267/banknote+authentication}

\begin{center}
\begin{tabular}{|l|l|}
\hline
Instances & 1372\\
\hline
Features & 4 Numerical Features\\
\hline
Classes & 2\\
\hline
Missing Values & None\\
\hline
Task & Binary Classification\\
\hline
\end{tabular}
\end{center}

\textbf{Feature Description}

\begin{itemize}
\item Variance
\item Skewness
\item Curtosis
\item Entropy
\end{itemize}

Target Class

\begin{itemize}
\item 0 -- Authentic Banknote
\item 1 -- Forged Banknote
\end{itemize}

%----------------------------------------------------------

\section*{5. Source Code}

The complete source code for this experiment, including data loading, preprocessing, the from-scratch Perceptron implementation, training, and evaluation, is available on GitHub:

\url{https://github.com/tonypanda7/PLACEHOLDER-REPO-NAME}

%----------------------------------------------------------

\section*{6. Mandatory Plots}

\begin{center}
\includegraphics[width=0.85\textwidth]{figures/histograms.png}
\end{center}

Variance and skewness show a broad, near-symmetric spread, whereas curtosis is right-skewed and entropy is concentrated toward negative values with a long left tail. This uneven spread across features indicates that feature scaling is necessary before training, since the perceptron's updates are sensitive to feature magnitude.

\vspace{3mm}

\begin{center}
\includegraphics[width=0.65\textwidth]{figures/correlation_heatmap.png}
\end{center}

Variance shows a positive correlation with skewness and a negative correlation with curtosis, while entropy is only weakly correlated with the other three features. This indicates some redundancy between variance, skewness and curtosis, but each feature still contributes distinct information that is useful for classification.

\vspace{3mm}

\begin{center}
\includegraphics[width=0.65\textwidth]{figures/scatter_plot.png}
\end{center}

The scatter plot of variance against skewness shows the two classes forming largely separate clusters, with authentic notes (class 0) occupying the higher-variance region and forged notes (class 1) clustering at lower variance values. A small region of overlap exists near the class boundary, but the data is close to linearly separable using just these two features.

\vspace{3mm}

\begin{center}
\includegraphics[width=0.85\textwidth]{figures/boxplots.png}
\end{center}

The boxplots confirm that variance and skewness have wider interquartile ranges and more symmetric distributions than curtosis and entropy. Curtosis and entropy display more points beyond the whiskers, suggesting these two features are noisier and less reliable individually for separating the two classes.

\vspace{3mm}

\begin{center}
\includegraphics[width=0.65\textwidth]{figures/training_error.png}
\end{center}

The training error drops sharply from 59 misclassified samples in epoch 1 to around 19--20 by epoch 5, after which it decreases slowly and plateaus near 17--18 errors for the remaining epochs. This shows that the training data is not perfectly linearly separable in the full 4-feature space, since the perceptron does not converge to zero error within 50 epochs.

\vspace{3mm}

\begin{center}
\includegraphics[width=0.65\textwidth]{figures/weight_evolution.png}
\end{center}

All weights change rapidly during the first few epochs and then stabilize, moving only slightly after around epoch 20. This gradual flattening is consistent with the training-error curve and indicates that the decision boundary is close to its final orientation well before training ends.

\vspace{3mm}

\begin{center}
\includegraphics[width=0.65\textwidth]{figures/bias_evolution.png}
\end{center}

The bias term becomes increasingly negative over the course of training, moving from 0 at initialization to about $-0.10$ by epoch 50. This shift indicates that the decision boundary is progressively translated to better separate the two classes as training proceeds.

\vspace{3mm}

\begin{center}
\includegraphics[width=0.65\textwidth]{figures/decision_boundary.png}
\end{center}

Using only variance and skewness, the perceptron learns a boundary that separates most authentic notes from forged notes, with a small number of points on the wrong side of the boundary near the overlap region. This confirms that these two features alone are highly informative, though not perfectly linearly separable, consistent with the near-perfect (but not 100\%) accuracy obtained with the full feature set.

\vspace{3mm}

\textbf{Note:} The Confusion Matrix and Learning Rate Comparison were evaluated numerically rather than as separate plots; the corresponding values are reported in Section 7 (Performance Tables).

%----------------------------------------------------------

\section*{7. Performance Tables}

\textbf{Training Summary}

\begin{tabular}{|p{6cm}|p{7cm}|}
\hline
Dataset Size & 1372 samples, 4 features, 2 classes\\
\hline
Train/Test Split & 1097 train (80\%) / 275 test (20\%), stratified\\
\hline
Learning Rate & 0.01\\
\hline
Epochs & 50 (run to completion; error did not reach zero)\\
\hline
Final Weights & $w_1=-0.1941$, $w_2=-0.2358$\\
\hline
Final Bias & $-0.1000$\\
\hline
Accuracy & 0.9855 (98.55\%)\\
\hline
Precision & 0.9683\\
\hline
Recall & 1.0000\\
\hline
F1-score & 0.9839\\
\hline
\end{tabular}

\vspace{3mm}

\textbf{Epoch-wise Learning}

\begin{tabular}{|c|c|c|c|c|}
\hline
Epoch & Errors & Weight 1 & Weight 2 & Bias\\
\hline
1 & 59 & -0.0680 & -0.0938 & -0.0300\\
2 & 25 & -0.0775 & -0.1042 & -0.0400\\
3 & 20 & -0.0802 & -0.1127 & -0.0400\\
4 & 19 & -0.0883 & -0.1196 & -0.0500\\
5 & 20 & -0.0991 & -0.1295 & -0.0500\\
10 & 20 & -0.1307 & -0.1541 & -0.0600\\
20 & 19 & -0.1636 & -0.1946 & -0.0700\\
30 & 17 & -0.1751 & -0.2043 & -0.0800\\
40 & 18 & -0.1948 & -0.2244 & -0.0900\\
50 & 18 & -0.1941 & -0.2358 & -0.1000\\
\hline
\end{tabular}

\vspace{3mm}

\textbf{Confusion Matrix (Test Set)}

\begin{tabular}{|p{4cm}|c|c|}
\hline
 & Predicted: Authentic (0) & Predicted: Forged (1)\\
\hline
Actual: Authentic (0) & 149 & 4\\
\hline
Actual: Forged (1) & 0 & 122\\
\hline
\end{tabular}

\vspace{3mm}

\textbf{Learning Rate Comparison}

\begin{tabular}{|c|c|c|}
\hline
Learning Rate & Final Train Errors (Epoch 50) & Test Accuracy\\
\hline
0.001 & 18 & 0.9855\\
0.01 & 18 & 0.9855\\
0.1 & 18 & 0.9855\\
\hline
\end{tabular}

\vspace{3mm}

\textbf{Comparison with Scikit-learn's Perceptron}

\begin{tabular}{|p{4cm}|c|c|}
\hline
Metric & From-Scratch Implementation & Scikit-learn Perceptron\\
\hline
Accuracy & 0.9855 & 0.9745\\
\hline
Precision & 0.9683 & 0.9457\\
\hline
Recall & 1.0000 & 1.0000\\
\hline
F1-score & 0.9839 & 0.9721\\
\hline
\end{tabular}

%----------------------------------------------------------

\section*{8. Discussion and Conclusion}

\subsection*{Discussion}

The perceptron achieved strong performance on the Banknote Authentication dataset (Accuracy 98.55\%, Precision 96.83\%, Recall 100\%, F1-score 98.39\%), consistent with the near-linear separability observed in the scatter and decision-boundary plots. The training error decreased sharply within the first five epochs and then plateaued near 17--18 misclassified samples, showing that the Perceptron Learning Rule does not fully converge because the training data is not perfectly linearly separable in the full 4-feature space; this is corroborated by the identical test accuracy obtained across learning rates 0.001, 0.01 and 0.1. Compared to Scikit-learn's Perceptron, which achieved a lower accuracy of 97.45\% and precision of 94.57\% under equivalent settings, the from-scratch implementation performed comparably or slightly better, validating the correctness of the manually implemented weighted sum, step activation, and weight/bias update rules. A recall of 1.0 with zero false negatives shows that no forged note was ever classified as authentic, while four authentic notes were misclassified as forged, which is the more acceptable direction of error for this application.

\subsection*{Conclusion}

This experiment successfully implemented a Single Layer Perceptron from scratch and used it to classify banknotes as authentic or forged, achieving 98.55\% test accuracy. The results confirm that the dataset's four features are strongly, though not perfectly, linearly separable. They also illustrate a key theoretical limitation of the Single Layer Perceptron: because the training data is not perfectly linearly separable, the perceptron learning rule does not converge to zero error and instead plateaus, motivating the use of multilayer architectures and differentiable activation functions for more complex, non-linearly separable problems.

%----------------------------------------------------------

\section*{9. References}

\begin{enumerate}
\item F. Rosenblatt, ``The Perceptron,'' \emph{Psychological Review}, 1958.
\item I. Goodfellow, Y. Bengio and A. Courville, \emph{Deep Learning}, MIT Press, 2016.
\item C. M. Bishop, \emph{Pattern Recognition and Machine Learning}, Springer, 2006.
\item S. Haykin, \emph{Neural Networks and Learning Machines}, Pearson, 2009.
\item UCI Machine Learning Repository -- Banknote Authentication Dataset.
\item Scikit-learn Documentation: Perceptron.
\end{enumerate}

\end{document}

